% Part: incompleteness
% Chapter: incompleteness-provability
% Section: rosser-thm

\documentclass[../../../include/open-logic-section]{subfiles}

\begin{document}

\olfileid[hi]{inc}{inp}{ros}

\olsection{रॉसर का प्रमेय}

क्या हम ग्योडेल (G\"odel) के प्रमाण को संशोधित करके अधिक प्रबल परिणाम पा
सकते हैं, जिसमें ``$\omega$-संगत'' के स्थान पर केवल ``संगत'' हो? रॉसर
(Rosser) की खोजी हुई एक युक्ति से उत्तर ``हाँ'' है। रॉसर की युक्ति
$\OProv[\Th{T}](y)$ के स्थान पर ``संशोधित'' !!{derivability} विधेय
$\ORProv_T(y)$ का उपयोग करना है।

\begin{thm}
\ollabel{thm:rosser}
$\Th{T}$,~$\Th{Q}$ का कोई संगत, !!{axiomatizable} विस्तार हो। तब
$\Th{T}$ पूर्ण नहीं है।
\end{thm}

\begin{proof}
याद करें कि $\OProv[\Th{T}](y)$ को
$\lexists[x][\OPrf[\Th{T}](x, y)]$ के रूप में परिभाषित किया जाता है, जहाँ
$\OPrf[\Th{T}](x, y)$ उस निर्णेय संबंध का निरूपण करता है जो तभी और केवल
तभी लागू होता है जब $x$, ग्योडेल (G\"odel) संख्या~$y$ वाले !!{sentence} की किसी
!!a{derivation} की ग्योडेल (G\"odel) संख्या हो। वह संबंध भी निर्णेय है जो $x$ और~$y$
के बीच तब लागू होता है जब $x$, ग्योडेल (G\"odel) संख्या~$y$ वाले वाक्य के किसी
\emph{खंडन} की ग्योडेल (G\"odel) संख्या हो। $\fn{not}(x)$ को वह आदिम पुनरावर्ती फलन
मानें जो यह करता है: यदि $x$ किसी सूत्र $!A$ का कूट है, तो
$\fn{not}(x)$,~$\lnot !A$ का कूट है। तब $\Refut[\Th{T}](x, y)$ तभी और केवल
तभी लागू होता है जब $\Prf[\Th{T}](x, \fn{not}(y))$। उसका निरूपण
$\ORefut[\Th{T}](x, y)$ करे। तब यदि $\Th{T} \Proves \lnot !A$ और
$\delta$ कोई संगत !!{derivation} है, तो
$\Th{Q} \Proves \ORefut[\Th{T}](\gn{\delta}, \gn{!A})$। हम
$\ORProv[\Th{T}](y)$ को इस प्रकार परिभाषित करते हैं:
\[
\lexists[x][(\OPrf[\Th{T}](x,y) \land \lforall[z][(z < x \lif \lnot
  \ORefut[\Th{T}](z,y))])].
\]
मोटे तौर पर $\ORProv[\Th{T}](y)$ कहता है, ``$\Th{T}$ में $y$ का कोई
प्रमाण है, और $y$ का इससे छोटा कोई खंडन नहीं है।'' यदि $\Th{T}$ संगत हो,
तो $\ORProv[\Th{T}](y)$ उन्हीं संख्याओं के लिए सत्य है जिनके लिए
$\OProv[\Th{T}](y)$ सत्य है; किंतु~$\Th{T}$ में \emph{सिद्धता} की दृष्टि
से (और अब हम जानते हैं कि सत्यता और सिद्धता में अंतर है!) दोनों के गुण
अलग हैं। यदि $\Th{T}$ \emph{अ}संगत है, तो दोनों समान संख्याओं के लिए लागू
\emph{नहीं} होते!{} ($\ORProv[\Th{T}](y)$ को प्रायः ``$y$ रॉसर-सिद्ध है''
पढ़ा जाता है। चूँकि अभी की चर्चा के अनुसार रॉसर-सिद्धता कोई विशेष प्रकार की
सिद्धता नहीं है---असंगत सिद्धांतों में ऐसे !!{sentence}s होते हैं जो सिद्ध
तो हैं, किंतु रॉसर-सिद्ध नहीं---यह भ्रमकारी हो सकता है। भ्रम से बचने के
लिए आप इसके स्थान पर ``$y$ श्मूवेबल है'' पढ़ सकते हैं।)

नियत-बिंदु उपपत्ति से कोई सूत्र $!R_\Th{T}$ ऐसा है कि
\begin{equation}
  \Th{Q} \Proves !R_\Th{T} \liff \lnot \ORProv[\Th{T}](\gn{!R_\Th{T}}).
  \ollabel{RT}
\end{equation}
\olref[1in]{thm:first-incompleteness} के प्रमाण के विपरीत, यहाँ हमारा दावा
है कि यदि $\Th{T}$ संगत है, तो $\Th{T}$ न $!R_\Th{T}$ को !!{derive} करता है,
और $\Th{T}$ न $\lnot !R_\Th{T}$ को !!{derive} करता है। (दूसरे शब्दों में हमें
$\omega$-संगतता की मान्यता नहीं चाहिए।)

पहले दिखाएँ कि $\Th{T} \Proves/ !R_{\Th{T}}$। मान लें वह इसे सिद्ध करता
है; तब~$T$ से~$!R_{\Th{T}}$ की कोई !!a{derivation} है। उसकी ग्योडेल (G\"odel) संख्या
$n$ मानें। तब
$\Th{Q} \Proves \OPrf[\Th{T}](\num{n}, \gn{!R_{\Th{T}}})$, क्योंकि
$\OPrf[\Th{T}]$,~$\Th{Q}$ में $\Prf[\Th{T}]$ का निरूपण करता है। इसके
अतिरिक्त प्रत्येक $k < n$ के लिए $k$,~$\lnot !R_{\Th{T}}$ की किसी
!!a{derivation} की ग्योडेल (G\"odel) संख्या नहीं है, क्योंकि $\Th{T}$ संगत है। अतः
प्रत्येक $k < n$ के लिए
$\Th{Q} \Proves \lnot \ORefut[\Th{T}](\num{k}, \gn{!R_{\Th{T}}})$।
\olref[req][min]{lem:less-nsucc} से
$\Th{Q} \Proves \lforall[z][(z < \num{n} \lif \lnot \ORefut[\Th{T}](z,
  \gn{!R_{\Th{T}}}))]$। अतः
\[
\Th{Q} \Proves \lexists[x][(\OPrf[\Th{T}](x,\gn{!R_{\Th{T}}}) \land \lforall[z][(z
    < x \lif \lnot \ORefut[\Th{T}](z,\gn{!R_{\Th{T}}}))])],
\]
किंतु यह केवल $\ORProv[\Th{T}](\gn{!R_{\Th{T}}})$ है। \olref{RT} से
$\Th{Q} \Proves \lnot !R_{\Th{T}}$। चूँकि $\Th{T}$,~$\Th{Q}$ का विस्तार
है, इसलिए $\Th{T} \Proves \lnot !R_{\Th{T}}$ भी। हमने मान लिया था कि
$\Th{T} \Proves !R_{\Th{T}}$, अतः $\Th{T}$ असंगत होगा, जो प्रमेय की
मान्यता के विरुद्ध है।

अब दिखाएँ कि $\Th{T} \Proves/ \lnot !R_{\Th{T}}$। फिर मान लें कि वह इसे
सिद्ध करता है और $n$,~$\lnot !R_{\Th{T}}$ की किसी !!a{derivation} की
ग्योडेल (G\"odel) संख्या है। तब $\Refut[\Th{T}](n, \Gn{!R_{\Th{T}}})$ लागू होता है;
और चूँकि $\ORefut[\Th{T}]$,~$\Th{Q}$ में $\Refut[\Th{T}]$ का निरूपण करता
है, इसलिए
$\Th{Q} \Proves \ORefut[\Th{T}](\num{n}, \gn{!R_{\Th{T}}})$। हम फिर
दिखाएँगे कि तब $\Th{T}$ असंगत होगा, क्योंकि वह~$!R_{\Th{T}}$ को भी
!!{derive} करेगा। चूँकि
\begin{align*}
\Th{Q} & \Proves !R_{\Th{T}} \liff \lnot \ORProv[\Th{T}](\gn{!R_{\Th{T}}}),
\intertext{और चूँकि $\Th{T}$,~$\Th{Q}$ का विस्तार है, इतना दिखाना पर्याप्त है कि}
\Th{Q} & \Proves \lnot \ORProv[\Th{T}](\gn{!R_{\Th{T}}}).
\end{align*}
!!{sentence} $\lnot \ORProv[\Th{T}](\gn{!R_{\Th{T}}})$, अर्थात्
\begin{align*}
  \lnot & \lexists[x][(\OPrf[\Th{T}](x,\gn{!R_{\Th{T}}}) \land
    \lforall[z][(z < x \lif \lnot \ORefut[\Th{T}](z,\gn{!R_{\Th{T}}}))])],
  \intertext{तार्किकतः इसके समतुल्य है:}
  & \lforall[x][(\OPrf[\Th{T}](x,\gn{!R_{\Th{T}}}) \lif
    \lexists[z][(z < x \land \ORefut[\Th{T}](z,\gn{!R_{\Th{T}}}))])].
\end{align*}
हम तर्क और $\Th{Q}$ द्वारा !!{derive} किए जाने वाले तथ्यों का उपयोग करके
अनौपचारिक तर्क देते हैं। मान लें $x$ स्वेच्छ है और
$\OPrf[\Th{T}](x, \gn{!R_{\Th{T}}})$। हम पहले ही जानते हैं कि
$\Th{T} \Proves/ !R_{\Th{T}}$; अतः प्रत्येक $k$ के लिए
$\Th{Q} \Proves \lnot \OPrf[\Th{T}](\num{k}, \gn{!R_{\Th{T}}})$। इसलिए
प्रत्येक $k$ के लिए $\eq/[x][\num{k}]$ निकलता है। विशेषतः हमारे पास
(क) $\eq/[x][\num{n}]$ है। हमारे पास
$\lnot(\eq[x][\num{0}] \lor \eq[x][\num{1}] \lor \dots \lor
\eq[x][\num{n-1}])$ भी है, और इसलिए
\olref[req][min]{lem:less-nsucc} से (ख) $\lnot(x < \num{n})$।
\olref[req][min]{lem:trichotomy} से $\num{n} < x$। चूँकि
$\Th{Q} \Proves \ORefut[\Th{T}](\num{n}, \gn{!R_{\Th{T}}})$, हमारे पास
$\num{n} < x \land \ORefut[\Th{T}](\num{n}, \gn{!R_{\Th{T}}})$ है, और
उससे $\lexists[z][(z < x \land
\ORefut[\Th{T}](z,\gn{!R_{\Th{T}}}))]$। चूँकि $x$ स्वेच्छ था, हमें
आवश्यकतानुसार मिलता है कि
\[
\lforall[x][(\OPrf[\Th{T}](x,\gn{!R_{\Th{T}}}) \lif
  \lexists[z][(z < x \land \ORefut[\Th{T}](z,\gn{!R_{\Th{T}}}))])].
\]
\end{proof}

\begin{prob}
प्राकृतिक संख्याओं के दो समुच्चय $A$ और $B$ \emph{संगणनीयतः अवियोज्य}
कहलाते हैं यदि ऐसा कोई !!{decidable} समुच्चय~$X$ न हो कि
$A \subseteq X$ और $B \subseteq \Complement{X}$
($\Complement{X}$,~$X$ का पूरक $\Nat \setminus X$ है)। $\Th{T}$ को
$\Th{Q}$ का कोई संगत !!{axiomatizable} विस्तार मानें। मान लें $A$ उन
!!{sentence}s की ग्योडेल (G\"odel) संख्याओं का समुच्चय है जो~$\Th{T}$ में सिद्ध हैं,
और $B$ उन वाक्यों की ग्योडेल (G\"odel) संख्याओं का समुच्चय है जो~$\Th{T}$ में खंडनीय
हैं। सिद्ध करें कि $A$ और~$B$ संगणनीयतः अवियोज्य हैं।
\end{prob}

\end{document}
